% 自动生成：work/gen_appendix_notes.py —— 由各 chunk 的 % [Tutor Note: ...] 汇集，勿手工编辑
\chapter*{附录 A\quad 译校批注}
\addcontentsline{toc}{chapter}{附录 A\quad 译校批注}
\markboth{附录 A\quad 译校批注}{附录 A\quad 译校批注}

本附录按章汇集正文中的译校批注，便于读者把批注与原文对照阅读。批注分三类，均在条目开头标明：\textbf{背景}给出原文未及展开的史实、人物与文献线索；\textbf{难点}说明推导中省略的中间步骤、约定与可能的误区；\textbf{校勘}记录与原文（或与通行文献）不一致之处及本译本的处理。批注编号中的\textbf{原书 p.$N$}表示该批注对应的日文原书页码（原书页码 $=$ PDF 页码 $-\,9$）。\par\medskip

\section*{A.1 第 1 章 序章}
\addcontentsline{toc}{section}{A.1 第 1 章 序章}
\begin{itemize}[leftmargin=2.4em,itemsep=4pt,label={\textbullet}]
\item \textbf{原书 p.1} 「アカデミア」即柏拉图于公元前 387 年前后在雅典西北郊创立的学园（Akadēmeia），门匾传说见诸后世记述，真伪不详，但常被引作几何学地位的象征。
\item \textbf{原书 p.1} 阿基米德墓上刻圆柱与内切球，典出 Cicero《Tusculanae disputationes》；他证明球体积为外切圆柱体积的 2/3，并希望以此作为自己的墓志。
\item \textbf{原书 p.2} 「ハレー彗星」的回归预言见 Halley 1705 年《Astronomiae cometicae synopsis》；1705 年他据 Newton 力学算出周期约 75 年，并预言 1758 年回归。原书此处以「Halley ら」称之。
\item \textbf{原书 p.3} 原书作「エルランゲン目録」，指 Klein 1872 年的 Erlangen Programm（埃尔朗根纲领）：以变换群对几何分类。
\item \textbf{原书 p.3} 「シンプレクティック」一词由 Weyl 在 1939 年《The Classical Groups》中引入，取自希腊语 sym-plektikos（交织、缠绕）。
\item \textbf{原书 p.3} 三大作图问题的不可解性分别由 Wantzel（1837，三等分角、倍积立方体）与 Lindemann（1882，化圆为方，依据 $\pi$ 的超越性）证明。
\item \textbf{原书 p.4} 「等位集合」指所有守恒量取定值的公共水平集；在 Liouville–Arnold 定理的假设下它是 $n$ 维环面 $T^n$，即 Lagrange 环面，见 6.1 节。
\item \textbf{原书 p.5} 原书此处作「古典力学から量子力学への以降」，「以降」当为「移行」之误，译文按「过渡」处理。
\end{itemize}

\section*{A.2 第 2 章 分析力学}
\addcontentsline{toc}{section}{A.2 第 2 章 分析力学}
\begin{itemize}[leftmargin=2.4em,itemsep=4pt,label={\textbullet}]
\item \textbf{原书 p.6} 背景：Newton 三定律见于《自然哲学的数学原理》(1687)；第一定律的实质是断言惯性系的存在，而非仅描述一种运动状态。
\item \textbf{原书 p.7} 背景：Noether 定理（1918）把连续对称性与守恒量对应起来：时间平移对应能量、空间平移对应动量、旋转对应角动量。
\item \textbf{原书 p.9} 背景：Kepler 问题是二体问题；解空间维数 12，对称性使其约化到质心匀速直线运动（3 维）+ 相对运动的平面轨道（2 维）。
\item \textbf{原书 p.9} 难点：(2.1.12) 与 (2.1.13) 中的力大小相同、方向相反，正是第三定律。
\item \textbf{原书 p.10} 背景：角动量守恒对应旋转对称性，这是 Kepler 第二定律（面积速度一定）的内蕴原因。
\item \textbf{原书 p.10} 难点：(2.1.25) 用了 $\dot{\bm{x}}\times\dot{\bm{x}}=0$ 与 (2.1.19)；$\bm{x}\times(-\bm{x}/|\bm{x}|^3)=0$ 因两向量共线。
\item \textbf{原书 p.11} 难点：(2.1.35) 右端 $\grad f\cdot\dot{\bm{x}}$ 即 $f$ 沿轨道的方向导数；$\partial f/\partial x$ 等是对坐标分量的偏导。
\item \textbf{原书 p.11} 校勘：原书此处作「(2.1.32) を」，由上下文（由 (2.1.33) 改写而来）应为 (2.1.33)，此处照录原文。
\item \textbf{原书 p.12} 校勘：由 (2.1.49) 有 $M^2=el$，故 (2.1.50) 末项应为 $el/r-1$、(2.1.51) 应为 $er\cos\theta=el-r$；原书印作 $el/4-1$、$el-1$，此处照录原文。
\item \textbf{原书 p.13} 背景：Fermat 原理严格地说应表述为光程取驻值（一般并非最小），原书此处作「最小」，译文照原书。
\item \textbf{原书 p.16} 校勘：例 2.2.3、例 2.2.4 原书 Lagrange 函数均印作 $L(x,v)$，按坐标应分别为 $L(\theta,\omega)$、$L(\bm{x}_1,\bm{x}_2,\bm{v}_1,\bm{v}_2)$；(2.2.25) 第二项按推导应含因子 $m$。均照录原文。
\item \textbf{原书 p.17} 难点：(2.2.31) 即二体问题分解为重心运动（质量 $m_1+m_2$、坐标 $\bm{x}'$）与相对运动（约化质量、坐标 $\bm{x}$）两部分。
\item \textbf{原书 p.18} 背景：取 $V=T_qW$、$f=L(q,\cdot)$，则 (2.3.4) 给出 Hamilton函数 $H(q,p)=p\cdot\dot q-L(q,\dot q)$，$p=\partial L/\partial\dot q$ 即共轭动量。
\item \textbf{原书 p.21} 难点：此处把 $\bm{p}$ 视为与 $\bm{q}$ 独立的变量（不再要求$\bm{p}=\partial L/\partial\dot{\bm{q}}$），这是相空间表述的关键一步。
\item \textbf{原书 p.22} 背景——$\bm{p}\cdot\rd\bm{q}$ 是余切丛 $T^*W$ 上的重言 1 形式（tautological form），本译注后文亦称 Liouville 形式；其外微分给出 $T^*W$ 上的标准辛形式。
\item \textbf{原书 p.23} 难点——(2.6.8) 中 $\bm{p}\cdot\dot{\bm{x}}=m\dot{x}^2/\sqrt{1-\dot{x}^2/c^2}$，故 $H=mc^2\bigl(1-\dot{x}^2/c^2+(\dot{x}/c)^2\bigr)/\sqrt{1-\dot{x}^2/c^2}=mc^2/\sqrt{1-\dot{x}^2/c^2}$。
\item \textbf{原书 p.23} 难点——(2.7.4) 由 (2.7.2) 代入 (2.7.3) 而得：第一项 $-\partial f/\partial p_i\cdot\partial H/\partial q_i$，第二项 $+\partial f/\partial q_i\cdot\partial H/\partial p_i$，两项之和恰为 (2.7.1) 中的 $\{H,f\}$。
\item \textbf{原书 p.24} 难点——(2.8.6)$\to$(2.8.7) 只用内层括号的反对称性 $\{f,h\}=-\{h,f\}$；(2.8.7)$\to$(2.8.8) 才是 Jacobi 恒等式。
\item \textbf{原书 p.25} 背景——Emmy Noether（1882--1935）1918 年证明的定理把连续对称性与守恒量对应起来；此处是其在哈密顿力学中的形式，(2.9.1) 即沿 $X_f$ 的方向导数。
\item \textbf{原书 p.26} 背景——$\theta$ 无理时每条叶在 $\T^2$ 中稠密，故 Kronecker 叶状结构是「叶非闭」的典型例子，与后面的辛叶（其叶为子流形）形成对照。
\item \textbf{原书 p.26} 校勘——(2.10.4) 原书作 $\Gamma(M,E)$，与 (2.10.3) 的 $\Gamma(E)$ 记号不一致，疑为 $\Gamma(E)$ 之误，此处照录。
\item \textbf{原书 p.27} 背景——Frobenius 定理（1877）是叶状结构理论的基础；条件 2 中 $\rd\mathcal{J}\subset\mathcal{J}$ 即「外微分不越出理想」，与对合性互为表里。
\item \textbf{原书 p.29} 背景——(2.11.9) 的 $\omega$ 正是辛叶上的辛形式，此命题说明「$\rd\omega=0$」与「Jacobi 恒等式」是同一条式子的两种说法。
\item \textbf{原书 p.29} 校勘——(2.11.17) 之后的正文原书作 $\omega(v,-)\colon T_x\to\R$，按上下文应为 $T_xM\to\R$，此处照录。
\item \textbf{原书 p.33} 校勘——(2.12.22) 原书作 $\bm{P}R(t)$，但由 (2.12.21) 与 $\bm{p}=\bm{P}R(-t)$，此处应为 $\bm{P}R(-t)$；译文照录原书。
\item \textbf{原书 p.35} 背景——Hill 区域得名于美国天文学家 G. W. Hill（1838--1914）的月球运动研究，其边界即零速度曲线。
\end{itemize}

\section*{A.3 第 3 章 G 结构的几何学}
\addcontentsline{toc}{section}{A.3 第 3 章 G 结构的几何学}
\begin{itemize}[leftmargin=2.4em,itemsep=4pt,label={\textbullet}]
\item \textbf{原书 p.36} 「Erlangen 大学的教授就任演讲」即 Klein 1872 年的埃尔朗根纲领（Erlanger Programm），其中提出以变换群对几何学分类的思想。
\item \textbf{原书 p.37} 由 Bieberbach 定理（1911–1912），紧致平坦流形必为环面被有限群自由作用所得之商。
\item \textbf{原书 p.38} Weyl 于 1939 年的《The classical groups》中自述此事；symplectic 与 complex 同源，本义为「编织在一起的」。
\item \textbf{原书 p.38} 原书用「包合的（involutive）」，即通称的余各向同性子空间（coisotropic subspace）；此处照原文译名。
\item \textbf{原书 p.39} 「Cartan 意义下的 G 结构」指切丛的结构群约化到 G 的结构；这里 Sp(2n,R) 结构即辛结构。
\item \textbf{原书 p.41} 此即 Moser 技巧：给定同调的变化路径，用非退化性对偶出向量场，再积分它得到所需的微分同胚；M 紧致保证向量场族可整体积分。
\item \textbf{原书 p.41} (3.2.16) 的求和上限原书作 2n，但 $p_i,q_i$ 只到 $n$，疑为 $n$ 之误，此处照录。
\item \textbf{原书 p.41} 校勘。原书此处写「$\R^n$ 中原点 $0$ 的邻域 $V$」；但下文 $\psi^*\omega_p$ 落在 $\bigwedge^2T_0^*\R^{2n}$ 中，辛形式也取 $\R^{2n}$ 的标准坐标，故原书此处的 $\R^n$ 应为 $\R^{2n}$。译文照录原文，读者请按 $\R^{2n}$ 理解。
\item \textbf{原书 p.42} 此处余切丛 $T^*L$ 赋以 4.1 节所述的标准辛结构；$L$ 作为零截面的像自然被视为 $T^*L$ 的 Lagrange 子流形。Weinstein 定理的证明与 Darboux 定理一样借助 Moser 定理，此处省略。
\item \textbf{原书 p.44} 对比定理 3.2.1：那里「$\omega$ 闭」是可积性条件，这里由 $N=0$ 取代 $\rd\omega=0$，故说 Nijenhuis 张量是 $\rd\omega$ 的复几何类比。
\item \textbf{原书 p.46} 此处 (3.5.13) 用了 $g$ 的 $J$ 不变性 $g(Y,JX)=g(JY,J(JX))$，(3.5.16) 用了 $g$ 的对称性。
\item \textbf{原书 p.47} 证明要点：把 $M$ 的基本 2 次形式限制到 $N$ 上即得 $N$ 的基本 2 次形式，而限制与外微分可交换，故闭性继承。
\item \textbf{原书 p.50} 此信条出自 Weinstein 1981 年的论文《Symplectic geometry》，意谓辛几何的对象与构造都可用 Lagrange 子流形重新表述。
\end{itemize}

\section*{A.4 第 4 章 各种辛流形}
\addcontentsline{toc}{section}{A.4 第 4 章 各种辛流形}
\begin{itemize}[leftmargin=2.4em,itemsep=4pt,label={\textbullet}]
\item \textbf{原书 p.52} 背景——完全辛流形即 $\omega$ 恰当的辛流形；余切丛上的 Liouville 向量场就是沿纤维方向伸缩的 Euler 向量场 $Z=\sum_i p_i\,\partial_{p_i}$，故称「Liouville」实与 Euler 齐次性同源。
\item \textbf{原书 p.52} 校勘——(4.2.2) 之前的正文原书作「$\C^k$ 的子簇」，但 $f_i$ 定义在 $\C^{2n+1}$ 上，同页下文亦作 $\C^{2n+1}$，疑为 $\C^{2n+1}$ 之误；照录原文。
\item \textbf{原书 p.52} 背景——Levi 问题问的是「多重劣调和函数的存在性是否刻画 Stein 流形」，由 Oka、Bremermann 等人在 1950 年代给出肯定答案（Oka 定理）。
\item \textbf{原书 p.53} 校勘——(4.3.3) 之后原书作「$G_\C$ 是 $\GL_{n+1}(\C)$ 或 $G=\SL_{n+1}(\C)$」，第二处记号当为 $G_\C=\SL_{n+1}(\C)$；照录原文。
\item \textbf{原书 p.55} 难点——(4.3.16)$\to$(4.3.18) 是对 $\log\sum|z_i|^2$ 逐次求 $\partial$、$\bar\partial$ 并把 $\bar\partial\partial\log$ 的分母通分的结果；三个式子都带同一个编号序列。
\item \textbf{原书 p.55} 校勘——(4.3.18) 原书在 $\bar z_i$ 与 $\rd z_i$ 之间排了 $\wedge$（即 $\bar z_i\wedge\rd z_i$），按 Fubini--Study 形式的通常写法此处应为函数与形式的乘积；照录原文。
\item \textbf{原书 p.56} 背景——$\mathfrak{g}^*$ 上的 Lie–Poisson 括号是 Kirillov、Kostant、Souriau 于 1960 年代建立的「轨道方法」的出发点；余伴随轨道上的辛形式即由该括号的辛叶结构诱导。
\item \textbf{原书 p.56} 难点——$Z=\nabla S$ 的验证：$\omega=\rd\rd^cS$，故 $\iota_{\nabla S}\omega = \rd S$ 的两边都在 $\rd\rd^cS$ 的框架下成立；关键是用 Kähler 度量把 1 形式 $\rd S$ 与向量场等同。
\item \textbf{原书 p.57} 背景——接触结构是「奇维的辛结构」：$N\times\R$ 上的辛化形式正是把 (4.6.2) 的 $\eta$ 沿 $t$ 方向延拓所得，故接触流形与辛流形的关系一如射影空间与仿射空间的关系。
\end{itemize}

\section*{A.5 第 5 章 动量映射与辛商}
\addcontentsline{toc}{section}{A.5 第 5 章 动量映射与辛商}
\begin{itemize}[leftmargin=2.4em,itemsep=4pt,label={\textbullet}]
\item \textbf{原书 p.58} 背景——把辛商与 GIT 商对应起来的是 Kempf--Ness 定理；改变 $\xi$ 即改变 GIT 中直线丛的线性化，其变化（VGIT）是 Thaddeus 等人研究的主题。
\item \textbf{原书 p.59} 难点——(5.1.2) 之所以能对每个 $Y$ 要求 $\underline{Y}$ 是 Hamilton 向量场，是因为 $G$ 保 $\omega$，故 $\iota_{\underline{Y}}\omega$ 是闭形式；在 $M$ 单连通时才自动恰当。
\item \textbf{原书 p.59} 校勘——(5.1.8) 的右端原书未加下划线，按 (5.1.2) 应理解为 $[\underline{Y},\underline{Z}]$；此处照录原文。
\item \textbf{原书 p.60} 难点——在 $\mathfrak{g}$ 的一组基上任意取常数再线性延拓，(5.1.4) 便成为线性映射；而 (5.1.13) 是额外要求，正是它定义了 Hamilton 作用而非一般的辛作用。
\item \textbf{原书 p.60} 难点——(5.1.15) 中 $e^{tY}\cdot x$ 关于 $t$ 的导数正是基本向量场 $\underline{Y}$ 在 $x$ 处的值，故对 $\mu_Z$ 求导即 $\underline{Y}\mu_Z$。
\item \textbf{原书 p.61} 背景——$\mathcal{O}$ 上的辛形式即 Kostant--Kirillov--Souriau 形式，它使 $G$ 的余伴随轨道成为齐性辛流形，与 4.4 节 $\mathfrak{g}^*$ 上的 Lie--Poisson 结构相容。
\item \textbf{原书 p.62} 难点——(5.1.34) 到 (5.1.37) 只是把 Cartan 同伦公式插入并整理；(5.1.40) 用到 $[\underline{Y},\underline{Z}]=\mathcal{L}_{\underline{Y}}\underline{Z}$，(5.1.41) 用到 $\mathcal{L}_{\underline{Y}}$ 是微分、与内乘的交换关系 $\mathcal{L}_{X}\iota_Y-\iota_Y\mathcal{L}_X=\iota_{[X,Y]}$。
\item \textbf{原书 p.64} 难点——(5.1.54) 是把 $d(e^{tY}\cdot z)/dt=(Ye^{tY})\cdot z$ 写成矩阵乘法后取 $t=0$；(5.1.61) 来自 $Y^*=-Y$ 即 $\overline{Y_{ij}}=-Y_{ji}$。
\item \textbf{原书 p.64} 背景——此即 Atiyah 与 Guillemin--Sternberg 的凸性定理（1982）；[10]、[55]、[57] 为相关原始文献，证明须用 Morse 理论，故原著从略。
\item \textbf{原书 p.65} 校勘——原书写「不依赖于 $q\in\pi^{-1}(p)$ 的取法」，按文意应为 $p\in\pi^{-1}(q)$（下文正是按 $p$ 在纤维中的选取来论证）；此处照录。
\item \textbf{原书 p.66} 校勘——(5.2.6) 之后原书把 $X,Y,X'',Y''$ 说成「$p$ 处 $\pi^{-1}(0)$ 的切向量」，但 $\pi$ 的定义域是 $\mu^{-1}(0)$（下文即用 $\pi$ 的商映射与 $\Ker\pi_*$），故 $\pi^{-1}(0)$ 应为 $\mu^{-1}(0)$，此处照录。
\item \textbf{原书 p.66} 难点——非退化性：$\omega$ 限制到 $\mu^{-1}(0)$ 后核的维数至多为 $\dim G$（由 (5.2.15)），而基本向量场张成的 $\dim G$ 维子丛已含于核中，故核恰为该子丛，商掉的正是它。
\item \textbf{原书 p.67} 背景——Marsden--Weinstein 约化由 Marsden 与 Weinstein 于 1974 年提出，Meyer 亦于 1973 年独立得到；对余伴随轨道取商把「动量取固定值」推广为「动量落在轨道内」。
\end{itemize}

\section*{A.6 第 6 章 完全可积系统}
\addcontentsline{toc}{section}{A.6 第 6 章 完全可积系统}
\begin{itemize}[leftmargin=2.4em,itemsep=4pt,label={\textbullet}]
\item \textbf{原书 p.68} 背景：角动量三分量之和不 Poisson 可交换，故不能同时取作 Darboux 坐标的一部分；可积性要求守恒量彼此可交换，这正是下面「函数独立 + Poisson 可交换」条件中交换性一项的来源。
\item \textbf{原书 p.69} 难点：(6.1.1) 的第二个等号用到 Leibniz 法则，并代入 $\{f_j,f_i\}=0$；因此复合函数不提供任何新的独立守恒量。
\item \textbf{原书 p.70} 背景：不变环面上的运动经作用–角变量线性化，这一「线性化」正是完全可积系统可解的含义；频率向量 $\omega=\partial H/\partial I$ 的算术性质决定轨道是周期还是稠密。
\item \textbf{原书 p.71} 校勘：原书 6.2.2 节标题与正文作「岩澤分解」（读作 Iwasawa），OCR 常误识为「岩津」；本译本统一作 Iwasawa 分解。
\item \textbf{原书 p.72} 背景：极分解 $g=ue^X$ 正是 $GL(n,\C)$ 的 Cartan 分解；$K=U(n)$ 与 Hermite 矩阵空间 $\mathfrak{p}$ 分别是极大紧子群与其正交补（对称空间 $GL(n,\C)/U(n)$ 的切空间）。
\item \textbf{原书 p.73} 难点：(6.2.12) 中的内积 $(v_i,u_j)$ 是 $\C^n$ 上的 Hermite 内积（对第二变量共轭线性）；因 $u_j$ 已正交归一，$w_i$ 与 $u_1,\dots,u_{i-1}$ 正交。
\item \textbf{原书 p.74} 背景：Iwasawa 分解 $G=KAN$ 是极分解的「非对称」加强版；在实矩阵情形 $K=O(n)$、$A$ 为正对角、$N$ 为单位上三角，恰是数值线性代数中的 QR 分解。
\item \textbf{原书 p.76} 校勘：原书此处作「$u\in B$」，据上下文（(6.2.38) 定义的 $u'$ 及其下标运算）应为「$u'\in B$」，此处照录并注明。
\item \textbf{原书 p.76} 校勘：原书 (6.2.20) 中 $w$ 取自 $\mathfrak{S}_n$，而此处起改用 Weyl 群 $W$；在 $GL(n,\C)$ 的情形二者一致（$W\cong\mathfrak{S}_n$），故不影响结论。
\item \textbf{原书 p.76} 校勘：原书作「Wely 群」，通行拼法为 Weyl 群（Hermann Weyl），此处照录并注明。
\item \textbf{原书 p.76} 背景：$G/B$ 即旗流形；Schubert 胞体 $Bw$ 的闭包给出 Schubert 簇，其维数等于置换 $w$ 的长度，这一胞体分解是 Schubert 演算与相交上同调的出发点。
\item \textbf{原书 p.77} 背景：Flaschka 因变量把非线性方程 (6.3.4) 化为 $a_i,b_i$ 的二次方程组，从而可以写成矩阵的 Lax 形式；这一步是 Toda 格子完全可积性的关键。
\item \textbf{原书 p.78} 难点：Lax 形式意味着 $L(t)$ 与 $L_0$ 相似，故特征值全都不随时间变化；Toda 格子的 $n$ 个独立守恒量即由 $L$ 的特征多项式给出。
\item \textbf{原书 p.79} 难点：(6.3.25) 的等式是把第一项写成 $-Q^{-1}\dot{Q}\cdot L$、第二项写成 $L\cdot(-Q^{-1}\dot{Q})$，其中用到 $Q^{-1}L_0\dot{Q}=Q^{-1}L_0QQ^{-1}\dot{Q}$。
\item \textbf{原书 p.80} 首行左端系数 4 来自 (6.4.1) 中的 1/2；右端于是化为 2[[RX,Y]+[X,RY],Z]\_R，逐项展开即得下列各式。
\item \textbf{原书 p.82} 末行照原书，求号只标 l,m 而 E\_i 亦须对 i 求和；此处不改动。
\item \textbf{原书 p.84} 量子 Yang--Baxter 方程 (6.4.33) 的 $\hbar$ 一次项给出经典方程，二次项给出 $s$ 的相容条件；此处只取一次项。
\item \textbf{原书 p.84} (6.5.1) 即 Lie--Poisson 括号在 $\mathfrak{g}^*$ 上的定义式 $\{H,\varphi\}(\xi)=\langle\xi,[(dH)_\xi,(d\varphi)_\xi]\rangle$。
\item \textbf{原书 p.85} AKS 取自 Adler、Kostant、Symes 三人姓氏的首字母，三人在 1970 年代各自独立得到该结论。
\item \textbf{原书 p.85} 原书作「包含性定理」，西文原注 involutivity；通行译名为「对合性」，此处从通行译法。
\item \textbf{原书 p.91} (6.7.19) 原书右端上界印作 $m-1$；按 (6.7.18) 与 $\xi\in f+\bigoplus_{i=-m}^{n-1}\mathfrak{g}^*_i$ 推来应为 $n-1$，此处照录原文。
\item \textbf{原书 p.94} (6.7.49) 原书右端对 $M_+$、$M_-$ 只印一个式子 $\frac{1}{2}(R+1)(d\varphi)_L$，结合 (6.7.52)(6.7.53) 应理解为该式的 $\pm$ 分量；此处照录原文。
\item \textbf{原书 p.95} 原书正文作「岩澤分解」；「岩澤」即 Iwasawa 的日文写法，故译作 Iwasawa 分解。
\end{itemize}

\section*{A.7 第 7 章 Hamilton 动力系统}
\addcontentsline{toc}{section}{A.7 第 7 章 Hamilton 动力系统}
\begin{itemize}[leftmargin=2.4em,itemsep=4pt,label={\textbullet}]
\item \textbf{原书 p.96} 背景：本章是「定性理论」一章，三条主线为 KAM（不变环面的存留）、Liouville（保测性）、遍历定理（时间平均=空间平均），与第 6 章的可积性理论相对照。
\item \textbf{原书 p.96} 难点：(7.1.2) 说明 $\Phi$ 是 $T$ 在 $M$ 上的作用；$T$ 取 $\N$ 得到频闪映射，取 $\R$ 得到流，二者是同一定义下的特例。
\item \textbf{原书 p.97} 背景：不变环面之所以特殊，是因为其上运动是准周期的、可以显式写出；KAM 理论正是研究这些环面在小摄动下是否存留。
\item \textbf{原书 p.98} 校勘：原书此处写 $k\in\Z^d\setminus\{0\}$，而振动数向量的维数为 $n$（见 (7.3.2)），$d$ 当为 $n$；此处照录原样。
\item \textbf{原书 p.98} 背景：「永年项」得名于天体力学：摄动解中随时间线性增长的项会长期（saeculum）累积，导致近似在长时间后失效。
\item \textbf{原书 p.99} 校勘：原书此句印作「正の Lebesgue 速度」，据上下文（补集的「測度」）应为「Lebesgue 測度」；此处译作测度并注明。
\item \textbf{原书 p.99} 背景：KAM 定理的结论只覆盖测度为正的 Cantor 集，共振区（补集）虽随 $\epsilon\to0$ 收缩，却在每一 $\epsilon$ 下仍可能存在混沌轨道。
\item \textbf{原书 p.99} 背景：Liouville 定理只要求保测，不要求可积；因此遍历论的一切工具都可施加于任意 Hamilton 系，这是第 7.5、7.6 节的前提。
\item \textbf{原书 p.99} 难点：(7.5.6) 到 (7.5.7) 的变形用 (7.5.5)：各层差集测度相等，故第二行每项都是 $\mu(E_1\setminus E_0)$；再与 $\mu(E_0)$ 相加得 $\mu(E)=\mu(E_0)$。
\item \textbf{原书 p.100} 背景：条件期望 $E(f\mid\mathfrak{G})$ 是「已知不变集信息后对 $f$ 的最佳猜测」；(7.6.2) 即其在 $L^2$ 意义下为 $f$ 向 $\mathfrak{G}$ 可测函数空间的正交投影。
\item \textbf{原书 p.101} 背景：遍历性即 $\mathfrak{G}$ 平凡，此时 $E(f\mid\mathfrak{G})=\int_X f\,d\mu$ 为常数，故 (7.6.3) 化为「时间平均=空间平均」，这正是统计力学的理论基础。
\end{itemize}

\section*{A.8 第 8 章 环面流形}
\addcontentsline{toc}{section}{A.8 第 8 章 环面流形}
\begin{itemize}[leftmargin=2.4em,itemsep=4pt,label={\textbullet}]
\item \textbf{原书 p.102} 「辛环面流形 $\leftrightarrow$ 凸多面体」的对应即下一节的 Delzant 定理；它把辛几何问题翻译成凸多面体的组合论问题，故常被称为辛几何的「离散化」。
\item \textbf{原书 p.102} 单值性的来源是 $GL_n(\Z)\ltimes\R^n$ 中仿射部分与线性部分的非平凡粘合：绕底空间中的回路一周回到原纤维时，作用坐标被整体平移、角坐标被 $GL_n(\Z)$ 整变换。
\item \textbf{原书 p.103} 「crepant 奇点消解」指消解映射 $f\colon\tilde X\to X$ 满足 $K_{\tilde X}=f^*K_X$（典范除子不改变）的奇点消解；三维普通二重点即 8.4.3 节 Atiyah flop 的出发点。
\item \textbf{原书 p.104} 原书此条行文作「多面体の $\Delta$ の単純かつ有理的な多面体 $\Delta$ は」，文句重复，疑为「多面体 $\Delta$ 是单纯且有理的，且满足下述条件时称为光滑」之讹；此处照其义译出。
\item \textbf{原书 p.104} 定理 8.1.3 是本章主干：它把「辛环面流形的同构分类」完全归结为「Delzant 多面体的同构分类」，即一个纯组合论问题；(2) 的构造在 p.105--106 给出。
\item \textbf{原书 p.105} (8.1.8) 中第二项 $\rd z\wedge\rd\bar z/|z|^2$ 恰是 $\C^\times$ 上不变形式 $\rd\log|z|\wedge\rd\arg z$ 的两倍，故其动量映射取 $\log|z|$；第一项对应 $\C$ 上的标准形式，动量映射为 $|z|^2/2$。
\item \textbf{原书 p.106} 此处辛商是在 $\mu_K$ 的「平移」水平集上取的；$\iota^*(-\lambda)$ 是 $\mathfrak{k}^*$ 中的正则值，Delzant 多面体的法向量数据保证 $K$ 作用自由，故商空间光滑。
\item \textbf{原书 p.106} 扇列与除子列的对偶关系是 Delzant 构造的枢纽：组合数据（多面体的法向量 $u_i$）经 $\varphi$ 变成环面同态，其核 $K$ 再经辛商还原出 $X$；除子列则说明 $M_\R$ 正是 $\R^d$ 中被「除掉」的 $K$ 方向的正交补。
\item \textbf{原书 p.107} 用 Moser 定理时，关键在于 $\omega_t$ 的上同调类恒定，故 $\dot\omega_t=\rd\sigma_t$ 有原像，解 Moser 方程 $\iota_{v_t}\omega_t=-\sigma_t$ 即得一族保持 $[\omega_t]$ 的微分同相。
\item \textbf{原书 p.109} 此处集合应为 $\bigcup_{x\in\Delta}\{N_x\Delta\}$；原书印作 $T_x\Delta$，与「法锥的集合」不合，属笔误。
\item \textbf{原书 p.110} 此处「$n\ge 1$ 个格点」应为「$k\ge 1$ 个格点」；原书后文随即出现「$n$ 与 $k$ 的常数」，可证此为笔误。
\item \textbf{原书 p.113} (8.4.32) 左端原作 $\C^2$，但此处动量映射定义在 $\C^3$ 上、商为 $\PP^2$，故应为 $\C^3$；原书笔误。
\item \textbf{原书 p.116} 「同義直線束」是原书用词，日文通常作「普遍直線束」，指 tautological bundle；此处按中文惯例译作重言线丛。
\item \textbf{原书 p.118} Atiyah flop 是三维极小模型理论中最早的分裂收缩例子：$X_+\to X_0$ 与 $X_-\to X_0$ 沿同一条 $\PP^1$ 收缩，把两个爆破互相替换，得到复流形同型但环面结构不同的两族辛环面流形。
\end{itemize}

\section*{A.9 第 9 章 二维规范理论}
\addcontentsline{toc}{section}{A.9 第 9 章 二维规范理论}
\begin{itemize}[leftmargin=2.4em,itemsep=4pt,label={\textbullet}]
\item \textbf{原书 p.119} Narasimhan–Seshadri 定理（1965）断言紧 Riemann 面上次数与秩互素的稳定全纯向量丛与基本群的中心扩张的不可约酉表示一一对应；Atiyah–Bott 将其重述为无限维辛商，见本章 9.5 节。
\item \textbf{原书 p.119} KdV 方程（Korteweg–de Vries）与 KP 方程（Kadomtsev–Petviashvili）都是可积的孤子方程；「可积层级」指由无穷多个彼此相容的方程构成的族。
\item \textbf{原书 p.120} 2 维时曲率张量只剩一个独立分量，故 (9.1.4) 右端只有单一项 $R$；这是「Einstein $\Longleftrightarrow$ 常曲率」在 2 维成立的原因。
\item \textbf{原书 p.120} Killing–Hopf 定理（1898）分类单连通完备常曲率空间：$S^n$、$\R^n$、$\HH^n$ 三种模型空间。
\item \textbf{原书 p.120} 「近复结构可积」指存在相容的复坐标图册，即 Newlander–Nirenberg 条件 $\bar\partial^2=0$ 成立；2 维时该条件自动满足。
\item \textbf{原书 p.121} 常中心曲率联络在文献中多称 Hermite–Einstein 联络；「中心」指曲率取值为李代数中心里的标量矩阵。
\item \textbf{原书 p.121} GAGA 是 Serre 论文标题「Géométrie Algébrique et Géométrie Analytique」的缩写，1956 年发表于 Ann. Inst. Fourier。
\item \textbf{原书 p.122} Atiyah–Bott（1982, Phil. Trans. Roy. Soc. A 308, 523–615）给出 9.5 节无限维辛商观点，是本章的理论源头。
\item \textbf{原书 p.123} (9.2.14) 的减号来自分次 Leibniz 律 $d(A\wedge s)=dA\wedge s+(-1)^{|A|}A\wedge \rd s$，此处 $|A|=1$。
\item \textbf{原书 p.124} (9.3.6) 用了 $X$ 上的内积 $\langle\alpha,\beta\rangle=\int\tr(\alpha\wedge*\beta)$ 的对称性与 $\tr$ 的循环性；(9.3.7) 用了一次分部积分并丢弃边界项。
\item \textbf{原书 p.125} $\mathrm{Gr}(r,N)$ 为 $\C^N$ 的 $r$ 维复子空间构成的 Grassmann 流形；$\PP^1=\mathrm{Gr}(1,2)$ 给出 $BU(r)$ 的 2 骨架，故 2 维底流形上只需一个整值不变量。
\item \textbf{原书 p.127} 原书此句引入 $\eta\in\Omega^1(P\times_G\mathfrak g)$，但下面 (9.5.21)–(9.5.23) 中只出现 $\alpha$；疑为笔误，照录。
\item \textbf{原书 p.130} 「与全纯结构相容」按定义即 $(0,1)$ 分量等于该全纯结构自带的 $\bar\partial_E$；这是后面标准联络存在唯一性的前提。
\item \textbf{原书 p.131} (9.7.7) 表明标准联络的 1 形式完全由 $\alpha=\bar\partial$ 的标架矩阵决定，故又称 Chern 联络。
\item \textbf{原书 p.133} (9.8.4) 右端的 $F^{0,2},F^{1,1}$ 原书未写下标 $A$，照录；实即 $F_A^{0,2},F_A^{1,1}$。
\item \textbf{原书 p.135} 原书把向量丛的「直既約」gloss 作 indecomposable，把联络的「直既約」gloss 作 irreducible，同一词形两种含义；译文分别作「直和不可分解」与「不可约」。
\item \textbf{原书 p.136} 上段「$v\in B$」原书如此，按上下文应为 $v\in V$（$B$ 疑为 $V$ 之误），照录。
\item \textbf{原书 p.137} 原书称 $L^p_k$ 范数被 $L^q_l$ 范数从上控制，方向似与嵌入 $L^p_k\subset L^q_l$ 相反（通常为 $\|f\|_{L^q_l}\le C\|f\|_{L^p_k}$）；照录原书。
\item \textbf{原书 p.140} (9.11.33) 右端原书作 $C^0(C)$，按上下文应为 $C^0(X)$（$C$ 疑为 $X$ 之误），照录。
\item \textbf{原书 p.141} 背景：$k-\frac np>0$ 保证 $L^p_k$ 中函数的乘积仍有 $k$ 阶弱导数，这就是 Sobolev 嵌入给出的乘法封闭性，故 $L^p_k$ 成为 Banach 代数。
\item \textbf{原书 p.141} 背景：原文本节及第 9.15 节把 $\mu(\cE)=\deg\cE/\rank\cE$ 称为「勾配」（本义为斜率、倾斜度）。依术语表统一译作「梯度」，全章同此；与 9.9 节标题「梯度与稳定性」一致。
\item \textbf{原书 p.142} 校勘：原书 (9.12.4) 把下排画成从右向左的正合列 $0\leftarrow\cQ_2\leftarrow\cE_2\leftarrow\cS_2\leftarrow0$，与上排方向相反；此处照录原图式，读时按 $0\to\cS_2\to\cE_2\to\cQ_2\to0$ 理解。
\item \textbf{原书 p.142} 背景：Dedekind 整域即代数曲线的坐标环（维数 1 的正则整域），其上有限生成无挠模必为局部自由模,这正是「层」与「丛」两套语言可以互换的代数根据（Serre 的 GAGA 型对应）。
\item \textbf{原书 p.142} 校勘：原书此处印作「$\iota_2\colon\cS_2\colon\cS_2\to\cE_2$」，多出一个「$\cS_2\colon$」,系排印重复，此处按其本意录作 $\iota_2\colon\cS_2\to\cE_2$。
\item \textbf{原书 p.142} 难点：设 $\varphi\colon M\to M'$ 为 $A$-模同态、$A$ 为整域。若 $m=fn$（$f\in K$、$n\in N=\ker\varphi$），则 $\varphi(m)=f\varphi(n)=0$，故 $m\in\ker\varphi$；即核在饱和下不变。
\item \textbf{原书 p.144} 背景：图 9.1 中以 $r_1,r_2$ 为宽、$d_1,d_2$ 为高作阶梯形折线与对角线，折线在直线上方当且仅当 $\mu(\cV_1)>\mu(\cV)$ 且 $\mu(\cV_2)<\mu(\cV)$。因倾斜方向如同跷跷板的两端，故称跷跷板原理。
\item \textbf{原书 p.145} 校勘：原书 (9.12.30) 前的不等式印作「$\mu(\cD)\ge\cD_i$」，右端脱一 $\mu$，当作 $\mu(\cD)\ge\mu(\cD_i)$；此处照原书录出。
\item \textbf{原书 p.150} 校勘：原书 p.150 行内两处下标疑倒置。其一印作「$\cF_i/\cF_{i+1}$ 与 $\cG_j/\cG_{j+1}$ 的稳定性」，因 $\cF_i\subset\cF_{i+1}$ 而应为 $\cF_{i+1}/\cF_i$、$\cG_{j+1}/\cG_j$；其二 (9.12.67) 分子第二项印作 $\cF_j$，按 Zassenhaus 引理 (9.12.55) 应为 $\cG_j$。此处照原书录出。
\item \textbf{原书 p.153} 校勘：原书 (9.13.19) 第一项求和下标印作 $i,j=1,\dots,n$，而 $g_{ij}$ 为 $m$ 次矩阵、区域在 $\R^m$ 中，当为 $i,j=1,\dots,m$；此处照原书录出。
\item \textbf{原书 p.155} 校勘：原书 (9.14.17) 内层印作「$|\lambda_1|e_{i1}|^2+\cdots+\lambda_r|e_{ir}|^2|$」，其中 $\lambda_r$ 外脱绝对值符号；因 $\lambda_j$ 为实特征值且外层取绝对值，结果不受影响，此处照原书录出。
\item \textbf{原书 p.157} 背景：原文以 $\mathcal{O}(\cE)$ 记 $\cE$ 上联络在复规范群 $\cG_{\C}$ 下的轨道（orbit），$J|_{\mathcal{O}(\cE)}$ 表示 $J$ 在该轨道上的限制；下文 9.15、9.16 各处的 $\mathcal{O}(\cE)$ 同此。
\item \textbf{原书 p.161} 背景：(9.15.16) 中 $\beta$ 只以一次幂出现于非对角块，故 $A(\cS,\cQ,t\beta)$ 当 $t\to0$ 时逐项收敛于分块对角的 $A(\cS,\cQ,0)$，各阶导数亦然，这就是「证明是显然的」的含义。
\item \textbf{原书 p.163} 难点：(9.15.46)(9.15.51) 把 $\nu$ 换成了 $\pm\tr$，依据是引理 9.14.4，须先确认括号内矩阵的特征值全为同号。这里 $\Lambda_{\cS}$、$\Lambda_{\cQ}$ 的对角元是逐次商的梯度，其符号由 (9.15.29)(9.15.30) 与稳定性的比较给出。
\item \textbf{原书 p.165} 背景：线丛情形只需看全纯同态 $\alpha$ 的零点。$\alpha$ 是 $\cE^\vee\otimes\cF$ 的全纯截面，故 $\deg\cF-\deg\cE$ 等于 $\alpha$ 的零点个数（计重数），它非负；等于 $0$ 当且仅当 $\alpha$ 处处非零、即 $\cE\cong\cF$。
\end{itemize}

\section*{A.10 第 10 章 全纯辛流形}
\addcontentsline{toc}{section}{A.10 第 10 章 全纯辛流形}
\begin{itemize}[leftmargin=2.4em,itemsep=4pt,label={\textbullet}]
\item \textbf{原书 p.167} 难点] $X$ 为 $2n$ 维时 $\sigma^n$ 是 $(n,n)$ 型的体积形式，故「$\sigma^n$ 处处非零」等价于 $\sigma$ 非退化。
\item \textbf{原书 p.168} 背景] 条件 $h^0(\Omega_X^p)=0$（$0<p<\dim X$）排除了环面与超凯勒因子的混合，使分解中的因子唯一到有限覆盖。
\item \textbf{原书 p.169} 难点] 把 $i^{-1}=-i$ 等代入 $-\frac12(x+ixi+jxj+kxk)$，四个 $i,j,k$ 项两两相消，即得 $a-bi-cj-dk$。
\item \textbf{原书 p.171} 难点] 该二重覆盖 $S^3\cong SU(2)\to SO(3)$ 即 $Spin(3)$，$\pi_1(SO(3))=\Z/2$ 即由此得到。
\item \textbf{原书 p.173} 难点] 超凯勒商是辛商的三重叠加：每个动量映射各降 $2\dim G$ 维，三个合起来恰好降 $4\dim G$ 维。
\item \textbf{原书 p.173} 难点] 写法 $z+jw$ 中的 $j$ 是不作右乘的虚单位，$z,w$ 取复坐标；故 $dz_i\wedge dw_i$ 为 $(2,0)$ 型，$\sigma_I$ 全纯。
\item \textbf{原书 p.174} 背景] 同一环面作用在 $\C^d$ 上取辛商得辛环面流形，在 $T^*\C^d$ 上取超凯勒商得超环面流形，后者是前者的超凯勒类比。
\item \textbf{原书 p.174} 背景] 原书节名为「箙多様体」；「箙」即 quiver（箭图），故本节标题按其数学内容译作「箭图簇」。
\item \textbf{原书 p.175} 难点] (10.7.6) 中的 $\varepsilon$ 取自 $\{\pm1\}$，其符号随复结构 $I,J,K$ 的选取而变，原书未在此处说明。
\item \textbf{原书 p.175} 背景] 箭图簇把几何表示论中的 Nakajima 对应与超凯勒商联系起来：取 $W_v=0$ 时所得即 ALE 空间上的瞬子模空间。
\end{itemize}

\section*{A.11 后记}
\addcontentsline{toc}{section}{A.11 后记}
\begin{itemize}[leftmargin=2.4em,itemsep=4pt,label={\textbullet}]
\item \textbf{原书 p.176} 背景。原书在此只提「邦訳」（日译本）；中译本读者更常用的是中译本：Landau–Lifshitz《力学》（理论物理学教程第一卷，高等教育出版社有译本），Arnold《经典力学的数学方法》（科学出版社／高等教育出版社有译本）。读者若找不到日译本，可径用中译本对照，公式编号一致。
\item \textbf{原书 p.176} 难点。Stokes 定理 $\int_M \rd\omega = \int_{\partial M}\omega$ 的「伴随」说法，指的是 $(d\omega, \eta) = (\omega, \delta\eta)$ 这一形式伴随关系：外微分 d 与余微分 $\delta$ 在 $L^2$ 内积下互为伴随，而 $\delta$ 正是边缘算子 $\partial$ 在链复形层面的对偶。原书此处说的是「边界作用素」与「外微分作用素」互为伴随作用素，宜理解为：链层面 $\partial$ 与上链层面 $d$ 的对偶配对，以及 Hodge 星号把两者联系起来。
\item \textbf{原书 p.176} 背景。森田茂之《微分形式の幾何学》与松本幸夫《多様体の基礎》都是日本流形论的标准入门书，前者以 de Rham 理论见长，后者以定义细致、图多著称。
\item \textbf{原书 p.176} 校勘。原书正文作「複素幾何」（复（素）几何）；中译名「复素几何」中的「素」为日文「複素＝complex」的直译，中文习惯作「复几何」。此处照录原书名，正文行文一律用「复几何」。
\item \textbf{原书 p.177} 校勘。原书此处作「第 5 章を書く際には」；但所列 [110]、[13]、[高崎 01] 均为完全可积系统（第 6 章）的文献，且上下文正讨论第 6 章，故疑为「第 6 章」之误。
\item \textbf{原书 p.179} 校勘。原书此句为「[三松 01] が日本数学会のホームページからダウンロードできる」，指该文（数理解析研究所講究録别册? 实为日本数学会「数学通信」／「季刊·数学」所载解说）可从学会网页下载 PDF，而非指书。
\item \textbf{原书 p.180} 校勘。原书此处作「同次直線束」，扫描件字迹与 OCR 均不清（OCR 读作「同音」）；「同次直線束」即齐性直线丛（homogeneous line bundle），$\PP^1$ 上的 $\OO(-1)$ 正是齐性直线丛，故译作「同次直线丛」；第 8 章等处译作「重言线丛」，二者同指。
\item \textbf{原书 p.180} 背景。(1)–(7) 是 conifold 的「形变」描述：$Y_\varepsilon$ 是以 $S^3$ 为底的 $\R^3$ 丛，即 $TS^3$，它与 $T^*S^3$ 辛同胚——这正是「$T^*S^3$ 的零截面在 $S^3$ 的 conifold 形变中」的确切含义。$\varepsilon\to0$ 时 $Y_\varepsilon$ 退化到奇点 $w_1^2+\cdots+w_4^2=0$。
\item \textbf{原书 p.181} 背景。佐佐木流形（Sasaki manifold）的 Riemann 锥为 Kähler 这一等价，是 Sasaki–Einstein 几何的出发点；$S^5$ 的锥就是 $\C^3\setminus\{0\}$，故 $S^5$ 与 $S^2\times S^3$ 的圆锥都对应 Kähler 锥，这正是 AdS/CFT 里把 $S^5$ 换成 Sasaki–Einstein 流形的几何依据。
\item \textbf{原书 p.181} 校勘。原书「超対称殻ゲージ理論」的「殻」字义不可通（OCR 亦不确定），按上下文（Klebanov–Witten 模型、quiver gauge theory）应读作「超对称规范理论」，或原文残缺一字。中译本按「超对称规范理论」处理。
\item \textbf{原书 p.183} 难点。(11) 中的 $\Lambda$ 是关于 Kähler 形式的内乘（与 $\omega$ 作缩并），$\Lambda F_A=\mu\,\id$ 是「常中心曲率」条件的矩阵形式：$F_A$ 的 $\omega$-迹（即中心部分）为常数 $\mu$ 乘以恒等。$X$ 为曲线时 $F_A$ 是 $(1,1)$ 形式，$\Lambda F_A$ 就是标量曲率，故与「常中心曲率」一致。
\item \textbf{原书 p.186} 术语。原书「箙多様体」即 quiver variety，中译通作「箭图簇」。此处与第 10 章译名统一。
\item \textbf{原书 p.186} 难点。(24) 中 $\Ext^1(\mathcal{E},\mathcal{E})^*$ 与 $T^*_{[\mathcal{E}]}\mathcal{M}$ 的等同，来自形变理论：$\Ext^1$ 分类 $\mathcal{E}$ 的一阶形变，故是切空间；对偶即余切空间。Higgs 场 $\Phi\in H^0(\omega_X\otimes\End\mathcal{E})$ 由 Serre 对偶正是余切向量，这解释了为何 Higgs 丛模空间 = 余切丛，从而自动带有超 Kähler 结构（余切丛总有 hyperkähler 度量）。
\end{itemize}
